\section{Split evaluations suffice without product preservation}

\subsection{F-SM42: the relative-injectivity criterion for the actual interval}

Retain precisely the ringoid, scalar normalization, and exact
resolution adjunction (SM40.1).  Thus
\(\mathcal M=\operatorname{Mod}\!-\!R\), the scalar
\(R(x,y)\) are projective, and the canonical scalar
derived counit at \(k\) is a quasi-isomorphism.
Let \(K\) be a \(k\)-linear endofunctor with a natural
map \(\alpha:\mathrm{Id}\to K\).  Suppose
\[
 \begin{gathered}
       L\varphi(KX)=0\quad(X\in\mathcal M),\\
       \alpha_X(x):X(x)\longrightarrow(KX)(x)
          \text{ is a split scalar monomorphism}
                    \quad(X\in\mathcal M,\ x\in R).
 \end{gathered}                                                \tag{SM42.1}
\]
The scalar retractions need not be natural in \(x\), or
even in \(X\).  No preservation of products, coproducts,
cokernels, or exact sequences is required of \(K\).
Then the ordinary counit \(\Lambda G(C)\to C\) is a
quasi-isomorphism for every connective scalar complex.
In particular the specified adjunction induces
\(\mathcal M[W^{-1}]\simeq D_{\ge0}(k)\), where
\(W=\Lambda^{-1}(\mathrm{qis})\), with its actual unit and counit.
Moreover, for arbitrary scalar modules \(U_i\),
\[
                  L\varphi\!\left(\prod_i C_{x_i}(U_i)\right)=0,
 \qquad C_x(U)(y)=\operatorname{Hom}_k(R(x,y),U).              \tag{SM42.2}
\]

\begin{proof}
Put \(I=\prod_iC_{x_i}(U_i)\).  The coinduction adjunction
identifies
\[
 \operatorname{Hom}_{\mathcal M}(X,I)
          =\prod_i\operatorname{Hom}_k(X(x_i),U_i).
                                                               \tag{SM42.3}
\]
The identity of \(I\) therefore corresponds to a family
of scalar maps \(q_i:I(x_i)\to U_i\).  Choose scalar
retractions \(r_i:(KI)(x_i)\to I(x_i)\) of
\(\alpha_I(x_i)\).  The maps \(q_i r_i\) correspond,
by the same adjunction, to an actual \(R\)-module map
\(r:KI\to I\).  Its composite with \(\alpha_I\)
corresponds to the original family \(q_i\), so
\(r\alpha_I=1_I\).  Thus \(I\) is a retract of
\(KI\), and (SM42.1) proves (SM42.2).

For an acyclic projective scalar chain object
\(P=\prod_{n\ge1}D_n(Q_n)\), the disk adjunction gives
\[
               GP=\prod_{n\ge1}\operatorname{Hom}_k(B_{n-1}(-),Q_n).
                                                               \tag{SM42.4}
\]
Each projective \(B_{n-1}\) is a retract of a coproduct
of left representables.  Hence \(GP\) is a retract of
a product of coinduced modules.  Formula (SM42.2) gives
\(L\varphi(GP)=0\).  The canonical scalar normalization
extends to free scalar modules and then projective
retracts, because \(\kappa\) and \(\Lambda\) are left
adjoints and scalar direct sums are exact.  Thus the
counit is a quasi-isomorphism on every projective chain
object, including its degree-zero projective summand.
The exact-adjunction criterion F-SM17, whose extension
uses a first-quadrant resolution with finite diagonals,
gives the result for every target and descends the actual
adjunction.  This is also precisely Step 5 of F-SM40;
its proof uses no further property of \(K\).
\end{proof}

\paragraph{Application to pointwise tensor products.}
Let \(\mathcal C\) be a small ordinary category,
\(R=k[\mathcal C]\), and
\(\mathcal M=\operatorname{Fun}(\mathcal C^{\mathrm{op}},
\operatorname{Mod}k)\), with pointwise tensor and
constant tensor unit \(\underline{k}\).  The scalar
functor is the exact constant functor; its left and
right adjoints are colimit and limit.  Retain the
canonical normalization and full interval acyclicity.
If the interval map
\(\alpha:\underline{k}\to J\) splits over \(k\) at
every object of \(\mathcal C\), then F-SM42 applies
to \(KX=J\otimes_kX\).  Indeed a scalar retraction
\(r_c:J(c)\to k\) gives
\(r_c\otimes1_{X(c)}\) as a retraction of
\(\alpha_X(c)\), for every \(X\).

For example, this scalar splitting holds whenever the
underlying interval is a reduced free presheaf
\[
             J(c)=k[I(c)]/k[s_0(c)],\qquad
                     \alpha_c(1)=[s_1(c)],                  \tag{SM42.5}
\]
where \(I\) has two natural sections \(s_0,s_1\) which
are distinct at every object.  Projecting onto the basis
element \([s_1(c)]\) gives the required scalar
retraction.  This projection is not asserted to be
natural in \(c\).  The bases \(I(c)\) can be infinite;
no finite morphism-set or finite interval hypothesis
has entered.  The normalization and full acyclicity
still have to be verified for a proposed interval;
the free-module description alone does not imply them.

\paragraph{The added range is strict: an actual interval which does not preserve products.}
For every nonzero commutative ring \(k\), take
\(\mathcal M=\operatorname{Ch}_{\ge0}(k)\), with its
usual homological tensor product and scalar unit
\(k[0]\).  In the standard ringoid presentation,
evaluation at a projective generator is evaluation
of a chain degree, and all scalar morphism modules
are zero or \(k\).  The scalar left adjoint is
\(\varphi=H_0\), with
\(L\varphi(X)\simeq X\) by the finite-diagonal
identification (SM39.11) for \(\mathcal N=\operatorname{Mod}k\).
Thus the canonical normalization holds.

Let \(V=k^{(\mathbb N)}\), with basis \(e_0,e_1,\ldots\),
put \(J=D_1(V)\), and take
\(\alpha:k[0]\to J\) to send \(1\) to \(e_0\)
in degree zero.  For every \(X\),
\[
  (J\otimes X)_n=(V\otimes X_n)\oplus(V\otimes X_{n-1}),
        \qquad \alpha_X(x)=(e_0\otimes x,0),                \tag{SM42.6}
\]
with \(X_{-1}=0\).  Projection onto the \(e_0\)
coefficient in the first summand splits this map over
\(k\) in each degree.  It is not claimed to be a chain
retraction.  If \(h:J_0\to J_1\) is the identity of
\(V\), extended by zero in the other degrees, then
\(dh+hd=1_J\).  The map \(h\otimes1_X\) contracts
\(J\otimes X\): the two cross terms involving \(d_X\)
cancel with their Koszul signs.  Therefore the full
interval acyclicity in (SM42.1) holds, including degree
zero.

Nevertheless \(K=J\otimes-\) does not preserve countable
products.  Apply its product comparison to the family
\(X_j=k[0]\), \(j\in\mathbb N\), and inspect degree
zero.  The comparison is
\[
 k^{(\mathbb N)}\otimes_k k^{\mathbb N}
                 \longrightarrow (k^{(\mathbb N)})^{\mathbb N}.
                                                               \tag{SM42.7}
\]
Every element on the left has nonzero coordinates in
only finitely many of the \(V\)-basis directions,
uniformly across all product factors.  The tuple
\((e_j)_{j\in\mathbb N}\) on the right has no such
finite support, and is not in the image.  This proves
failure of product preservation.  Thus F-SM42 genuinely
weakens F-SM40's hypothesis on the actual interval
functor.  This illustrative category was already
covered by F-SM39 and is not counted as a newly solved
example of the original question.

\paragraph{Residual scope.}
The split-evaluation condition is a sufficient condition,
not a converse assertion about the raw counit.  A case
still outside these proofs must have no verified interval
satisfying (SM42.1), in addition to escaping the earlier
scalar, resolution-finiteness, and connective-chain
conditions.  No derivation of scalar splitting from the
bare handwritten bilinearity convention is claimed.
